\documentclass[12pt,letterpaper]{article}
\usepackage[margin=15mm,top=22mm,bottom=18mm,headheight=14pt,headsep=10pt]{geometry}
\usepackage[T1]{fontenc}
\usepackage{helvet}
\renewcommand{\familydefault}{\sfdefault}
\usepackage{tikz}
\usetikzlibrary{arrows.meta}
\usepackage{fancyhdr}
\pagestyle{fancy}
\fancyhf{}
\newcommand{\band}{Grades 2--5}
\fancyhead[L]{\small Week 59 / Constant width / \band}
\fancyfoot[L]{\footnotesize Bellingham Math Circle / Week 59 / W59-S-v2}
\fancyfoot[R]{\footnotesize\thepage}
\renewcommand{\headrulewidth}{0pt}
\renewcommand{\footrulewidth}{0pt}
\setlength{\parindent}{0pt}
\setlength{\parskip}{7pt}
\newcommand{\problem}[2]{\textbf{Problem #1:} #2\par}
\newcommand{\spacebox}[1]{\vspace{3mm}\begin{tikzpicture}[x=1mm,y=1mm]\draw[gray!45,line width=.4pt] (0,0) rectangle (180,#1);\end{tikzpicture}\par}
\input{assets.tex}
\begin{document}
\fontsize{12.5}{15.5}\selectfont
To measure width, keep the two straight edges parallel, enclosing the whole piece and just touching it. Turn and slide the piece freely. Measure the perpendicular gap, to the nearest millimeter. This gap is the width in that position.

\measureExample

\problem{1}{Use the five pieces. Find positions that make each gap as small and as large as you can. Keep a sketch and the gap for each.}
\extremeTable

\newpage
\problem{2}{Keep the parallel edges fixed 60 mm apart while the piece turns and slides. Which pieces can turn all the way around inside them? Which can keep touching both edges throughout the turn? Keep a sketch of a deciding position.}
\turnTable

\newpage
\renewcommand{\band}{Grades 3--5}
A compass draws an arc with one center and one radius.

\compassExample

\problem{3}{Use the 60 mm and 40 mm equilateral templates to make two curved triangles. Each short arc joins two corners and has the third corner as its center. Predict the width of each piece, then test it.}
\constructionRecords
\spacebox{40}

\newpage
A line that touches a circle at just one point meets the radius there at a right angle.

\tangentExample

\problem{4}{Explain why the exact curved triangle has the same width in every direction. Its arcs have radius 60 mm, with centers at $A$, $B$ and $C$. Include the turns between the drawings in your explanation.}
\proofPositions
\spacebox{83}

\newpage
The mark $O$ is a point on a piece. Measure to one straight edge that encloses the whole piece and touches it.

\pointSupportExample

\problem{5}{Use the marked disk and curved triangle. Find the smallest and largest perpendicular distances from $O$ to one touching edge. Can a piece have constant width while those distances change?}
\heightTable
\spacebox{45}

\newpage
\renewcommand{\band}{Grades 4--5}
\problem{6}{The large circle has six equally spaced boundary dots. Compare the boundary lengths of the 60 mm-width disk and curved triangle. Give an exact argument.}
\perimeterCompare
\spacebox{105}
\end{document}
